\documentclass[11pt]{article}
% ============ FangMath self-contained preamble (pdflatex) ============
\usepackage[T1]{fontenc}
\usepackage{mathpazo}            % Palatino text + math
\usepackage[margin=0.9in,top=1in,bottom=1in]{geometry}
\usepackage{amsmath,amssymb}
\usepackage{xcolor}
\usepackage[most]{tcolorbox}
\usepackage{fancyhdr}
\usepackage{enumitem}
\usepackage{pifont}
\usepackage{needspace}
\usepackage{multicol}
\usepackage{tikz}
\usetikzlibrary{arrows.meta}

\definecolor{FangBlue}{HTML}{1F4E79}
\definecolor{FangLight}{HTML}{EAF1F8}
\definecolor{FangGreen}{HTML}{1D9E75}
\definecolor{FangRed}{HTML}{C0392B}
\definecolor{FangGold}{HTML}{B7862C}
\definecolor{FangPurple}{HTML}{6C3483}

% ---- Headers / footers ----
\setlength{\headheight}{26pt}
\pagestyle{fancy}
\fancyhf{}
\renewcommand{\headrulewidth}{0.8pt}
\renewcommand{\headrule}{\hbox to\headwidth{\color{FangBlue}\leaders\hrule height \headrulewidth\hfill}}
\lhead{\color{FangBlue}\textbf{FangMath}}
\rhead{\color{FangBlue}Absolute Value Inequalities --- Practice}
\cfoot{\color{FangBlue}\thepage}

% ---- Step box (from the In/Out notes) ----
\newtcolorbox{stepbox}[1]{
  colback=FangLight, colframe=FangBlue, boxrule=0.8pt, arc=3pt,
  left=10pt, right=10pt, top=6pt, bottom=6pt,
  fonttitle=\bfseries\color{white},
  colbacktitle=FangBlue, coltitle=white,
  title={#1}, breakable}
% ---- Mistake box ----
\newtcolorbox{mistakebox}{
  colback=white, colframe=FangRed, boxrule=1pt, arc=3pt,
  left=10pt, right=10pt, top=6pt, bottom=6pt,
  fonttitle=\bfseries\color{white},
  colbacktitle=FangRed, coltitle=white,
  title={Common Mistake to Avoid}, breakable}
% ---- Answer box ----
\newtcolorbox{answerbox}{
  colback=white, colframe=FangGreen, boxrule=1.2pt, arc=3pt,
  left=10pt, right=10pt, top=6pt, bottom=6pt}
\newcommand{\In}{\textbf{\color{FangGreen}in}}
\newcommand{\Out}{\textbf{\color{FangRed}out}}

% ---- Leveled practice box (structure borrowed from the equations worksheet) ----
\newtcolorbox{partbox}[2]{
  enhanced, breakable,
  colback=#1!7, colframe=#1, coltitle=white,
  fonttitle=\bfseries\large,
  title={#2},
  arc=3pt, boxrule=1.2pt,
  left=10pt, right=10pt, top=7pt, bottom=7pt}

% ---- Fill-in-the-blank helpers ----
\newcommand{\blank}{\underline{\hspace{2.6cm}}}
\newcommand{\sblank}{\underline{\hspace{1.3cm}}}
\newcommand{\ans}[1]{{\color{FangGreen}\bfseries #1}}
\newcommand{\wspA}{\vspace{0.3cm}}
\newcommand{\wspB}{\vspace{0.3cm}}
\newcommand{\wspC}{\vspace{0.4cm}}

% ---- Blank scratch number line for practice problems (fixed width, no numbers --
%      student plots boundaries/labels/shading themselves) ----
\newcommand{\blanknumline}{%
\begin{center}
\begin{tikzpicture}[x=0.5cm]
  \draw[thick,-{Stealth[length=2mm]}] (-6,0) -- (6,0);
  \draw[thick,-{Stealth[length=2mm]}] (6,0) -- (-6,0);
  \foreach \i in {-6,...,6} {\draw (\i,0.07) -- (\i,-0.07);}
\end{tikzpicture}
\end{center}}

% ---- Number line: just the two boundary dots, no shading yet ----
\newcommand{\numlineDots}[6]{%
\begin{center}
\begin{tikzpicture}[x=0.42cm]
  \draw[thick,-{Stealth[length=2mm]}] (#1-0.5,0) -- (#2+0.5,0);
  \draw[thick,-{Stealth[length=2mm]}] (#2+0.5,0) -- (#1-0.5,0);
  \filldraw[FangBlue] (#3,0) circle (3.2pt);
  \filldraw[FangBlue] (#4,0) circle (3.2pt);
  \node[below=5pt] at (#3,0) {#5};
  \node[below=5pt] at (#4,0) {#6};
\end{tikzpicture}
\end{center}}
% ---- Number line: shaded interval BETWEEN two closed endpoints (And / In) ----
\newcommand{\numlineIn}[6]{%
\begin{center}
\begin{tikzpicture}[x=0.42cm]
  \draw[thick,-{Stealth[length=2mm]}] (#1-0.5,0) -- (#2+0.5,0);
  \draw[thick,-{Stealth[length=2mm]}] (#2+0.5,0) -- (#1-0.5,0);
  \draw[line width=3pt,FangGreen] (#3,0) -- (#4,0);
  \filldraw[FangGreen] (#3,0) circle (3.2pt);
  \filldraw[FangGreen] (#4,0) circle (3.2pt);
  \node[below=5pt] at (#3,0) {#5};
  \node[below=5pt] at (#4,0) {#6};
\end{tikzpicture}
\end{center}}
% ---- Number line: shaded rays OUTSIDE two closed endpoints (Or / Out) ----
\newcommand{\numlineOut}[6]{%
\begin{center}
\begin{tikzpicture}[x=0.42cm]
  \draw[thick,-{Stealth[length=2mm]}] (#1-0.8,0) -- (#2+0.8,0);
  \draw[thick,-{Stealth[length=2mm]}] (#2+0.8,0) -- (#1-0.8,0);
  \draw[line width=3pt,FangGreen] (#1-0.8,0) -- (#3,0);
  \draw[line width=3pt,FangGreen] (#4,0) -- (#2+0.8,0);
  \filldraw[FangGreen] (#3,0) circle (3.2pt);
  \filldraw[FangGreen] (#4,0) circle (3.2pt);
  \node[below=5pt] at (#3,0) {#5};
  \node[below=5pt] at (#4,0) {#6};
\end{tikzpicture}
\end{center}}

% ---- Compact answer-key graphs (small scale, args: min max leftEndpoint rightEndpoint) ----
\newcommand{\gIn}[4]{%
\begin{center}
\begin{tikzpicture}[x=0.32cm]
  \draw[thick,-{Stealth[length=1.8mm]}] (#1-0.5,0) -- (#2+0.5,0);
  \draw[thick,-{Stealth[length=1.8mm]}] (#2+0.5,0) -- (#1-0.5,0);
  \draw[line width=2.4pt,FangGreen] (#3,0) -- (#4,0);
  \filldraw[FangGreen] (#3,0) circle (2.6pt);
  \filldraw[FangGreen] (#4,0) circle (2.6pt);
  \node[below=3pt,font=\footnotesize] at (#3,0) {$#3$};
  \node[below=3pt,font=\footnotesize] at (#4,0) {$#4$};
\end{tikzpicture}
\end{center}}
\newcommand{\gInOpen}[4]{%
\begin{center}
\begin{tikzpicture}[x=0.32cm]
  \draw[thick,-{Stealth[length=1.8mm]}] (#1-0.5,0) -- (#2+0.5,0);
  \draw[thick,-{Stealth[length=1.8mm]}] (#2+0.5,0) -- (#1-0.5,0);
  \draw[line width=2.4pt,FangGreen] (#3,0) -- (#4,0);
  \draw[FangGreen,thick,fill=white] (#3,0) circle (2.6pt);
  \draw[FangGreen,thick,fill=white] (#4,0) circle (2.6pt);
  \node[below=3pt,font=\footnotesize] at (#3,0) {$#3$};
  \node[below=3pt,font=\footnotesize] at (#4,0) {$#4$};
\end{tikzpicture}
\end{center}}
\newcommand{\gOut}[4]{%
\begin{center}
\begin{tikzpicture}[x=0.32cm]
  \draw[thick,-{Stealth[length=1.8mm]}] (#1-0.8,0) -- (#2+0.8,0);
  \draw[thick,-{Stealth[length=1.8mm]}] (#2+0.8,0) -- (#1-0.8,0);
  \draw[line width=2.4pt,FangGreen] (#1-0.8,0) -- (#3,0);
  \draw[line width=2.4pt,FangGreen] (#4,0) -- (#2+0.8,0);
  \filldraw[FangGreen] (#3,0) circle (2.6pt);
  \filldraw[FangGreen] (#4,0) circle (2.6pt);
  \node[below=3pt,font=\footnotesize] at (#3,0) {$#3$};
  \node[below=3pt,font=\footnotesize] at (#4,0) {$#4$};
\end{tikzpicture}
\end{center}}
\newcommand{\gOutOpen}[4]{%
\begin{center}
\begin{tikzpicture}[x=0.32cm]
  \draw[thick,-{Stealth[length=1.8mm]}] (#1-0.8,0) -- (#2+0.8,0);
  \draw[thick,-{Stealth[length=1.8mm]}] (#2+0.8,0) -- (#1-0.8,0);
  \draw[line width=2.4pt,FangGreen] (#1-0.8,0) -- (#3,0);
  \draw[line width=2.4pt,FangGreen] (#4,0) -- (#2+0.8,0);
  \draw[FangGreen,thick,fill=white] (#3,0) circle (2.6pt);
  \draw[FangGreen,thick,fill=white] (#4,0) circle (2.6pt);
  \node[below=3pt,font=\footnotesize] at (#3,0) {$#3$};
  \node[below=3pt,font=\footnotesize] at (#4,0) {$#4$};
\end{tikzpicture}
\end{center}}
\newcommand{\gPoint}[3]{%
\begin{center}
\begin{tikzpicture}[x=0.32cm]
  \draw[thick,-{Stealth[length=1.8mm]}] (#1-0.5,0) -- (#2+0.5,0);
  \draw[thick,-{Stealth[length=1.8mm]}] (#2+0.5,0) -- (#1-0.5,0);
  \filldraw[FangGreen] (#3,0) circle (2.8pt);
  \node[below=3pt,font=\footnotesize] at (#3,0) {$#3$};
\end{tikzpicture}
\end{center}}
\newcommand{\gAll}[2]{%
\begin{center}
\begin{tikzpicture}[x=0.32cm]
  \draw[line width=2.4pt,FangGreen,{Stealth[length=2.2mm]}-{Stealth[length=2.2mm]}] (#1-0.5,0) -- (#2+0.5,0);
  \node[below=3pt,font=\footnotesize] at (0,0) {all real numbers};
\end{tikzpicture}
\end{center}}

\begin{document}

% ─── TITLE BANNER ─────────────────────────────────────────────────────
\begin{tcolorbox}[enhanced, arc=6pt,
  colback=FangBlue, colframe=FangGreen, boxrule=2.2pt,
  left=12pt, right=12pt, top=10pt, bottom=10pt]
  \begin{center}
    {\color{white}\fontsize{19}{23}\selectfont\bfseries
      Solving Absolute Value Inequalities}\\[4pt]
    {\color{FangLight}\normalsize The ``In or Out'' Method --- Step-by-Step Practice}
  \end{center}
\end{tcolorbox}

\vspace{6pt}
\begin{tcolorbox}[colback=FangLight, colframe=FangBlue!50, arc=3pt, boxrule=0.8pt,
  left=8pt, right=8pt, top=5pt, bottom=5pt]
\begin{tabular}{p{6cm} p{4cm} p{4cm}}
Name:\;\underline{\hspace{5cm}} &
Date:\;\underline{\hspace{2.8cm}} &
Period:\;\underline{\hspace{2cm}}
\end{tabular}
\end{tcolorbox}

\vspace{8pt}
\begin{stepbox}{The Big Rule}
Absolute value bars ask: how far is the inside expression from $0$? Instead of memorizing which symbol
means \emph{and} and which means \emph{or}, \textbf{test the middle point} (where the inside expression
equals $0$) to see whether it belongs to the solution:
\begin{center}
\renewcommand{\arraystretch}{1.3}
\begin{tabular}{|c|c|c|}
\hline
\textbf{Test result at the middle} & \textbf{Region} & \textbf{Compound inequality} \\\hline
True & \In\ (between the boundaries) & \textbf{and} \\\hline
False & \Out\ (outside the boundaries) & \textbf{or} \\\hline
\end{tabular}
\end{center}
Every practice problem below uses the \textbf{same three steps}: (1) find the boundaries, (2) test the
middle point True/False, (3) write \emph{and}/\emph{or} and solve. Fill in every blank!
\end{stepbox}

\vspace{6pt}
\noindent{\color{FangBlue}\Large\bfseries Example 1 \quad Solve $\lvert 3x-4\rvert \le 8$ \normalsize (an ``And'' case)}
\vspace{2pt}

\begin{stepbox}{Step 1 \ Find the boundaries}
Since $\lvert -8\rvert = 8$ and $\lvert 8\rvert = 8$, the expression $3x-4$ hits the edge at
$-8$ and $8$. Mark both as \textbf{closed} dots (the inequality is $\le$, so the ends are included).
\numlineDots{-8}{8}{-8}{8}{$-8$}{$8$}
\end{stepbox}
\begin{stepbox}{Step 2 \ Check: In or Out?}
Test the middle point $3x-4 = 0$:
\[
\lvert 0\rvert = 0 \le 8 \quad\text{\textbf{\color{FangGreen}True}.}
\]
The point \emph{between} the boundaries works, so we keep the \In\ region. That makes this an
\textbf{and} statement (a \emph{between} inequality), not an \emph{or}.
\end{stepbox}
\begin{stepbox}{Step 3 \ Write the compound inequality and solve}
Because it is ``in,'' $3x-4$ stays \emph{between} $-8$ and $8$:
\[
-8 \le 3x-4 \le 8.
\]
Solve all three parts at once. Add $4$:
\[
-8+4 \le 3x-4+4 \le 8+4
\qquad\Longrightarrow\qquad
-4 \le 3x \le 12.
\]
Divide by $3$:
\[
\frac{-4}{3} \le x \le \frac{12}{3}
\qquad\Longrightarrow\qquad
-\frac{4}{3} \le x \le 4.
\]
\end{stepbox}
\begin{center}
\begin{answerbox}
\centering
\textbf{Solution:} \quad $\displaystyle -\frac{4}{3} \le x \le 4$
\quad or \quad $\left[-\dfrac{4}{3},\, 4\right]$
\end{answerbox}
\end{center}
\numlineIn{-2}{5}{-1.333}{4}{$-\frac{4}{3}$}{$4$}
\vspace{2pt}
\begin{center}
\small\itshape Check: at $x=0$, $\lvert 3(0)-4\rvert = 4 \le 8$ \checkmark
\end{center}

\vspace{10pt}
\noindent{\color{FangBlue}\Large\bfseries Example 2 \quad Solve $\lvert 2x+4\rvert \ge 6$ \normalsize (an ``Or'' case)}
\vspace{2pt}

\noindent Same method as Example 1 -- but this time the test point will fail, which flips it to an
\emph{or} statement.

\begin{stepbox}{Step 1 \ Find the boundaries}
Since $\lvert -6\rvert = 6$ and $\lvert 6\rvert = 6$, the expression $2x+4$ hits the edge at
$-6$ and $6$. Mark both as \textbf{closed} dots (the inequality is $\ge$, so the ends are included).
\numlineDots{-6}{6}{-6}{6}{$-6$}{$6$}
\end{stepbox}
\begin{stepbox}{Step 2 \ Check: In or Out?}
Test the middle point $2x+4 = 0$:
\[
\lvert 0\rvert = 0 \ge 6 \quad\text{\textbf{\color{FangRed}False}.}
\]
The point \emph{between} the boundaries does \textbf{not} work, so the solution is the \Out\ region.
That makes this an \textbf{or} statement (an \emph{outside} inequality), not an \emph{and}.
\end{stepbox}
\begin{stepbox}{Step 3 \ Write the compound inequality and solve}
Because it is ``out,'' $2x+4$ stays \emph{outside} $-6$ and $6$ --- less than or equal to the left
boundary, \textbf{or} greater than or equal to the right boundary:
\[
2x+4 \le -6 \qquad\text{\textbf{or}}\qquad 2x+4 \ge 6.
\]
Solve each piece on its own. Subtract $4$ from both sides:
\[
2x \le -10 \qquad\text{\textbf{or}}\qquad 2x \ge 2.
\]
Divide both sides by $2$:
\[
x \le -5 \qquad\text{\textbf{or}}\qquad x \ge 1.
\]
\end{stepbox}

\needspace{9\baselineskip}
\begin{mistakebox}
The original inequality is $\ge$, not a strict $>$ -- so the boundary points themselves are part of the
solution, and the compound inequality must use $\le$ and $\ge$, not $<$ and $>$. Writing ``$2x+4<-6$ or
$2x+4>6$'' would quietly kick $x=-5$ and $x=1$ back out of the answer. Always check: plug the boundary
value into the \emph{original} inequality --- at $x=-5$, $\lvert 2(-5)+4\rvert = \lvert -6\rvert = 6 \ge 6$
\checkmark, so it belongs, and the dots on the number line stay \textbf{closed}.
\end{mistakebox}

\begin{center}
\begin{answerbox}
\centering
\textbf{Solution:} \quad $x \le -5 \ \text{ or } \ x \ge 1$
\quad or \quad $(-\infty,-5\,] \cup [\,1,\infty)$
\end{answerbox}
\end{center}
\numlineOut{-9}{4}{-5}{1}{$-5$}{$1$}
\vspace{2pt}
\begin{center}
\small\itshape Check: at $x=-10$ (outside), $\lvert 2(-10)+4\rvert = 16 \ge 6$ \checkmark \qquad
at $x=0$ (between), $\lvert 2(0)+4\rvert = 4 \ge 6$? \ding{55}\ (correctly excluded)
\end{center}

\vspace{6pt}
\begin{stepbox}{Key Idea --- And vs.\ Or at a Glance}
\begin{center}
\renewcommand{\arraystretch}{1.3}
\begin{tabular}{|c|c|c|}
\hline
\textbf{Form} & \textbf{Test at the middle} & \textbf{Solution shape}\\\hline
$\lvert \text{expr}\rvert \le k$ & always True (In) & \textbf{and}: one connected interval\\\hline
$\lvert \text{expr}\rvert \ge k$ & always False (Out) & \textbf{or}: two separate rays\\\hline
\end{tabular}
\end{center}
(Assuming $k>0$.) The middle-point test is a shortcut for the graph of $y=\lvert\text{expr}\rvert$: it
dips \emph{below} the line $y=k$ only in the middle (giving an \emph{and}), so anywhere it stays at or
above $y=k$ gives an \emph{or}.
\end{stepbox}

\newpage
%% ══════════════════════════════════════════════════════════════════
%%  PART A --- Level 1 : direct |x|
%% ══════════════════════════════════════════════════════════════════
\begin{partbox}{FangBlue}{Part A \quad Level 1 --- Direct Absolute Value
  \hfill{\normalfont\small $|x| \; \square \; c$}}
\textit{Follow the three steps. The boundaries of $|x|$ are just $x=c$ and $x=-c$.}
\vspace{6pt}
\begin{multicols}{2}
\begin{enumerate}[label=\textbf{\arabic*.},itemsep=0pt]

\item $|x| \le 4$\\[3pt]
Step 1 --- Boundaries: $x=\sblank$ and $x=\sblank$\\[3pt]
Step 2 --- Test $x=0$: $|0|\;\underline{\hspace{0.7cm}}\;4$ $\Rightarrow$ True / False (circle one)\\[3pt]
In / Out \quad $\Rightarrow$ \quad and / or (circle one)\\[3pt]
Step 3: \sblank $\le x \le$ \sblank\\[3pt]
\textbf{Answer:} \blank\\[3pt]
Step 4 --- Graph it: \blanknumline
\wspB

\item $|x| \ge 3$\\[3pt]
Step 1 --- Boundaries: $x=\sblank$ and $x=\sblank$\\[3pt]
Step 2 --- Test $x=0$: $|0|\;\underline{\hspace{0.7cm}}\;3$ $\Rightarrow$ True / False (circle one)\\[3pt]
In / Out \quad $\Rightarrow$ \quad and / or (circle one)\\[3pt]
Step 3: $x\le$ \sblank\ \ or\ \ $x\ge$ \sblank\\[3pt]
\textbf{Answer:} \blank\\[3pt]
Step 4 --- Graph it: \blanknumline
\wspB

\item $|x| < 6$\\[3pt]
Step 1 --- Boundaries: $x=\sblank$ and $x=\sblank$\\[3pt]
Step 2 --- Test $x=0$: $|0|\;\underline{\hspace{0.7cm}}\;6$ $\Rightarrow$ True / False (circle one)\\[3pt]
In / Out \quad $\Rightarrow$ \quad and / or (circle one)\\[3pt]
Step 3: \sblank $< x <$ \sblank\\[3pt]
\textbf{Answer:} \blank\\[3pt]
Step 4 --- Graph it: \blanknumline
\wspB

\item $|x| > 2$\\[3pt]
Step 1 --- Boundaries: $x=\sblank$ and $x=\sblank$\\[3pt]
Step 2 --- Test $x=0$: $|0|\;\underline{\hspace{0.7cm}}\;2$ $\Rightarrow$ True / False (circle one)\\[3pt]
In / Out \quad $\Rightarrow$ \quad and / or (circle one)\\[3pt]
Step 3: $x<$ \sblank\ \ or\ \ $x>$ \sblank\\[3pt]
\textbf{Answer:} \blank\\[3pt]
Step 4 --- Graph it: \blanknumline
\wspB

\item $|x| \le 0$\\[3pt]
\textit{(Hint: absolute value can never be negative --- when can it equal exactly $0$?)}\\[3pt]
\textbf{Answer:} \blank\\[3pt]
Graph it: \blanknumline
\wspA

\item $|x| \ge 0$\\[3pt]
\textit{(Hint: absolute value is always $\ge 0$ --- for which $x$ is this ever false?)}\\[3pt]
\textbf{Answer:} \blank\\[3pt]
Graph it: \blanknumline
\wspA

\end{enumerate}
\end{multicols}
\end{partbox}

\vspace{10pt}
%% ══════════════════════════════════════════════════════════════════
%%  PART B --- Level 2 : shifted inside
%% ══════════════════════════════════════════════════════════════════
\begin{partbox}{FangPurple}{Part B \quad Level 2 --- Shift Inside
  \hfill{\normalfont\small $|x \pm a| \; \square \; c$}}
\textit{Same three steps --- but the middle point is no longer $x=0$; solve $x\pm a = 0$ first.}
\vspace{6pt}
\begin{multicols}{2}
\begin{enumerate}[label=\textbf{\arabic*.},itemsep=0pt]

\item $|x-2| \le 5$\\[3pt]
Step 1: $x-2=\sblank$ and $x-2=\sblank$ $\Rightarrow$ $x=\sblank$ and $x=\sblank$\\[3pt]
Step 2 --- Test middle ($x=2$): $|0|\;\underline{\hspace{0.7cm}}\;5$ $\Rightarrow$ True/False\\[3pt]
In/Out $\Rightarrow$ and/or (circle)\\[3pt]
Step 3: \sblank $\le x \le$ \sblank\\[3pt]
\textbf{Answer:} \blank\\[3pt]
Step 4 --- Graph it: \blanknumline
\wspC

\item $|x+4| \ge 3$\\[3pt]
Step 1: $x+4=\sblank$ and $x+4=\sblank$ $\Rightarrow$ $x=\sblank$ and $x=\sblank$\\[3pt]
Step 2 --- Test middle ($x=-4$): $|0|\;\underline{\hspace{0.7cm}}\;3$ $\Rightarrow$ True/False\\[3pt]
In/Out $\Rightarrow$ and/or (circle)\\[3pt]
Step 3: $x\le$ \sblank\ \ or\ \ $x\ge$ \sblank\\[3pt]
\textbf{Answer:} \blank\\[3pt]
Step 4 --- Graph it: \blanknumline
\wspC

\item $|x-5| < 2$\\[3pt]
Step 1: $x-5=\sblank$ and $x-5=\sblank$ $\Rightarrow$ $x=\sblank$ and $x=\sblank$\\[3pt]
Step 2 --- Test middle ($x=5$): $|0|\;\underline{\hspace{0.7cm}}\;2$ $\Rightarrow$ True/False\\[3pt]
In/Out $\Rightarrow$ and/or (circle)\\[3pt]
Step 3: \sblank $< x <$ \sblank\\[3pt]
\textbf{Answer:} \blank\\[3pt]
Step 4 --- Graph it: \blanknumline
\wspC

\item $|x+1| > 4$\\[3pt]
Step 1: $x+1=\sblank$ and $x+1=\sblank$ $\Rightarrow$ $x=\sblank$ and $x=\sblank$\\[3pt]
Step 2 --- Test middle ($x=-1$): $|0|\;\underline{\hspace{0.7cm}}\;4$ $\Rightarrow$ True/False\\[3pt]
In/Out $\Rightarrow$ and/or (circle)\\[3pt]
Step 3: $x<$ \sblank\ \ or\ \ $x>$ \sblank\\[3pt]
\textbf{Answer:} \blank\\[3pt]
Step 4 --- Graph it: \blanknumline
\wspC

\item $|x+3| \le 6$\\[3pt]
Step 1: $x+3=\sblank$ and $x+3=\sblank$ $\Rightarrow$ $x=\sblank$ and $x=\sblank$\\[3pt]
Step 2 --- Test middle ($x=-3$): $|0|\;\underline{\hspace{0.7cm}}\;6$ $\Rightarrow$ True/False\\[3pt]
In/Out $\Rightarrow$ and/or (circle)\\[3pt]
Step 3: \sblank $\le x \le$ \sblank\\[3pt]
\textbf{Answer:} \blank\\[3pt]
Step 4 --- Graph it: \blanknumline
\wspC

\item $|x-7| \ge 5$\\[3pt]
Step 1: $x-7=\sblank$ and $x-7=\sblank$ $\Rightarrow$ $x=\sblank$ and $x=\sblank$\\[3pt]
Step 2 --- Test middle ($x=7$): $|0|\;\underline{\hspace{0.7cm}}\;5$ $\Rightarrow$ True/False\\[3pt]
In/Out $\Rightarrow$ and/or (circle)\\[3pt]
Step 3: $x\le$ \sblank\ \ or\ \ $x\ge$ \sblank\\[3pt]
\textbf{Answer:} \blank\\[3pt]
Step 4 --- Graph it: \blanknumline
\wspC

\end{enumerate}
\end{multicols}
\end{partbox}

\vspace{10pt}
%% ══════════════════════════════════════════════════════════════════
%%  PART C --- Level 3 : isolate first
%% ══════════════════════════════════════════════════════════════════
\begin{partbox}{FangGold}{Part C \quad Level 3 --- Isolate First!
  \hfill{\normalfont\small Get $|\cdot|$ alone, then run the 3 steps}}
\textit{Add, subtract, or divide to isolate the absolute value first --- just like Example 2 required
solving for $2x+4$ before testing.}
\vspace{6pt}
\begin{enumerate}[label=\textbf{\arabic*.},itemsep=6pt]

\item $|x|+3 \le 7$\\[3pt]
Isolate: $|x| \le$ \sblank\\[3pt]
Step 1 --- Boundaries: $x=\sblank$ and $x=\sblank$\\[3pt]
Step 2 --- Test $x=0$: True / False\\[3pt]
In / Out $\Rightarrow$ and / or (circle)\\[3pt]
Step 3: \sblank $\le x \le$ \sblank\\[3pt]
\textbf{Answer:} \blank\\[3pt]
Graph it: \blanknumline

\item $2|x| \ge 10$\\[3pt]
Isolate: $|x| \ge$ \sblank\\[3pt]
Step 1 --- Boundaries: $x=\sblank$ and $x=\sblank$\\[3pt]
Step 2 --- Test $x=0$: True / False\\[3pt]
In / Out $\Rightarrow$ and / or (circle)\\[3pt]
Step 3: $x\le$ \sblank\ \ or\ \ $x\ge$ \sblank\\[3pt]
\textbf{Answer:} \blank\\[3pt]
Graph it: \blanknumline

\item $|x-1|+2 < 6$\\[3pt]
Isolate: $|x-1| <$ \sblank\\[3pt]
Step 1 --- Boundaries: $x-1=\sblank$ and $x-1=\sblank$ $\Rightarrow$ $x=\sblank$, $x=\sblank$\\[3pt]
Step 2 --- Test $x=1$: True / False \quad In / Out $\Rightarrow$ and / or\\[3pt]
Step 3: \sblank $< x <$ \sblank\\[3pt]
\textbf{Answer:} \blank\\[3pt]
Graph it: \blanknumline

\item $3|x+2| > 9$\\[3pt]
Isolate: $|x+2| >$ \sblank\\[3pt]
Step 1 --- Boundaries: $x+2=\sblank$ and $x+2=\sblank$ $\Rightarrow$ $x=\sblank$, $x=\sblank$\\[3pt]
Step 2 --- Test $x=-2$: True / False \quad In / Out $\Rightarrow$ and / or\\[3pt]
Step 3: $x<$ \sblank\ \ or\ \ $x>$ \sblank\\[3pt]
\textbf{Answer:} \blank\\[3pt]
Graph it: \blanknumline

\item $|x+4|-1 \le 3$\\[3pt]
Isolate: $|x+4| \le$ \sblank\\[3pt]
Step 1 --- Boundaries: $x+4=\sblank$ and $x+4=\sblank$ $\Rightarrow$ $x=\sblank$, $x=\sblank$\\[3pt]
Step 2 --- Test $x=-4$: True / False \quad In / Out $\Rightarrow$ and / or\\[3pt]
Step 3: \sblank $\le x \le$ \sblank\\[3pt]
\textbf{Answer:} \blank\\[3pt]
Graph it: \blanknumline

\item $4|x-3| \ge 8$\\[3pt]
Isolate: $|x-3| \ge$ \sblank\\[3pt]
Step 1 --- Boundaries: $x-3=\sblank$ and $x-3=\sblank$ $\Rightarrow$ $x=\sblank$, $x=\sblank$\\[3pt]
Step 2 --- Test $x=3$: True / False \quad In / Out $\Rightarrow$ and / or\\[3pt]
Step 3: $x\le$ \sblank\ \ or\ \ $x\ge$ \sblank\\[3pt]
\textbf{Answer:} \blank\\[3pt]
Graph it: \blanknumline

\end{enumerate}
\end{partbox}

\vspace{10pt}
%% ══════════════════════════════════════════════════════════════════
%%  PART D --- Challenge
%% ══════════════════════════════════════════════════════════════════
\begin{partbox}{FangRed}{Part D \quad Challenge
  \hfill{\normalfont\small Show all 3 steps on your own --- no scaffold}}
\vspace{4pt}
\begin{enumerate}[label=\textbf{\arabic*.},itemsep=14pt]

\item $|2x-4| \le 6$

\vspace{0.15cm}
Step 1 (boundaries): \underline{\hspace{5.5cm}}\\[10pt]
Step 2 (test the middle, True/False, In/Out, and/or): \underline{\hspace{5.5cm}}\\[10pt]
Step 3 (compound inequality, solved): \underline{\hspace{5.5cm}}

\vspace{0.2cm}
\textbf{Answer:} \blank\\[4pt]
Step 4 --- Graph it: \blanknumline

\item $|3x+3| \ge 9$

\vspace{0.15cm}
Step 1 (boundaries): \underline{\hspace{5.5cm}}\\[10pt]
Step 2 (test the middle, True/False, In/Out, and/or): \underline{\hspace{5.5cm}}\\[10pt]
Step 3 (compound inequality, solved): \underline{\hspace{5.5cm}}

\vspace{0.2cm}
\textbf{Answer:} \blank\\[4pt]
Step 4 --- Graph it: \blanknumline

\item $|2x+1| < 5$

\vspace{0.15cm}
Step 1 (boundaries): \underline{\hspace{5.5cm}}\\[10pt]
Step 2 (test the middle, True/False, In/Out, and/or): \underline{\hspace{5.5cm}}\\[10pt]
Step 3 (compound inequality, solved): \underline{\hspace{5.5cm}}

\vspace{0.2cm}
\textbf{Answer:} \blank\\[4pt]
Step 4 --- Graph it: \blanknumline

\item $|4x-8| \ge 0$
\hfill{\small\itshape Hint: can an absolute value ever be negative?}

\vspace{0.3cm}
\textbf{Answer:} \blank\\[4pt]
Graph it: \blanknumline

\end{enumerate}
\end{partbox}

%% ══════════════════════════════════════════════════════════════════
%%  ANSWER KEY
%% ══════════════════════════════════════════════════════════════════
\newpage
\begin{tcolorbox}[enhanced, breakable, arc=5pt,
  colback=FangGreen!10, colframe=FangGreen, coltitle=white,
  fonttitle=\bfseries\large,
  title={\quad Answer Key \quad (Teacher Copy)},
  boxrule=1.5pt, left=10pt, right=10pt, top=8pt, bottom=8pt]

\textbf{\color{FangBlue}Part A} --- each: boundaries, center test, and/or, final answer, graph
\begin{multicols}{2}
\begin{enumerate}[label=\textbf{\arabic*.},itemsep=4pt]
  \item $x=-4,4$; True; \ans{In, and} $\Rightarrow$ \ans{$-4\le x\le 4$}\\ \gIn{-6}{6}{-4}{4}
  \item $x=-3,3$; False; \ans{Out, or} $\Rightarrow$ \ans{$x\le -3$ or $x\ge 3$}\\ \gOut{-5}{5}{-3}{3}
  \item $x=-6,6$; True; \ans{In, and} $\Rightarrow$ \ans{$-6< x< 6$} (open circles)\\ \gInOpen{-8}{8}{-6}{6}
  \item $x=-2,2$; False; \ans{Out, or} $\Rightarrow$ \ans{$x< -2$ or $x> 2$} (open circles)\\ \gOutOpen{-4}{4}{-2}{2}
  \item \ans{$x=0$ (single point)}\\ \gPoint{-3}{3}{0}
  \item \ans{All real numbers}\\ \gAll{-4}{4}
\end{enumerate}
\end{multicols}

\vspace{4pt}\textbf{\color{FangPurple}Part B}
\begin{multicols}{2}
\begin{enumerate}[label=\textbf{\arabic*.},itemsep=4pt]
  \item $x=-3,7$; True; \ans{In, and} $\Rightarrow$ \ans{$-3\le x\le 7$}\\ \gIn{-5}{9}{-3}{7}
  \item $x=-7,-1$; False; \ans{Out, or} $\Rightarrow$ \ans{$x\le -7$ or $x\ge -1$}\\ \gOut{-9}{1}{-7}{-1}
  \item $x=3,7$; True; \ans{In, and} $\Rightarrow$ \ans{$3< x< 7$} (open circles)\\ \gInOpen{1}{9}{3}{7}
  \item $x=-5,3$; False; \ans{Out, or} $\Rightarrow$ \ans{$x< -5$ or $x> 3$} (open circles)\\ \gOutOpen{-7}{5}{-5}{3}
  \item $x=-9,3$; True; \ans{In, and} $\Rightarrow$ \ans{$-9\le x\le 3$}\\ \gIn{-11}{5}{-9}{3}
  \item $x=2,12$; False; \ans{Out, or} $\Rightarrow$ \ans{$x\le 2$ or $x\ge 12$}\\ \gOut{0}{14}{2}{12}
\end{enumerate}
\end{multicols}

\vspace{4pt}\textbf{\color{FangGold}Part C}
\begin{multicols}{2}
\begin{enumerate}[label=\textbf{\arabic*.},itemsep=4pt]
  \item $|x|\le 4$; $x=-4,4$; True; \ans{In, and} $\Rightarrow$ \ans{$-4\le x\le 4$}\\ \gIn{-6}{6}{-4}{4}
  \item $|x|\ge 5$; $x=-5,5$; False; \ans{Out, or} $\Rightarrow$ \ans{$x\le -5$ or $x\ge 5$}\\ \gOut{-7}{7}{-5}{5}
  \item $|x-1|<4$; $x=-3,5$; True; \ans{In, and} $\Rightarrow$ \ans{$-3< x< 5$} (open circles)\\ \gInOpen{-5}{7}{-3}{5}
  \item $|x+2|>3$; $x=-5,1$; False; \ans{Out, or} $\Rightarrow$ \ans{$x< -5$ or $x> 1$} (open circles)\\ \gOutOpen{-7}{3}{-5}{1}
  \item $|x+4|\le 4$; $x=-8,0$; True; \ans{In, and} $\Rightarrow$ \ans{$-8\le x\le 0$}\\ \gIn{-10}{2}{-8}{0}
  \item $|x-3|\ge 2$; $x=1,5$; False; \ans{Out, or} $\Rightarrow$ \ans{$x\le 1$ or $x\ge 5$}\\ \gOut{-1}{7}{1}{5}
\end{enumerate}
\end{multicols}

\vspace{4pt}\textbf{\color{FangRed}Part D}
\begin{multicols}{2}
\begin{enumerate}[label=\textbf{\arabic*.},itemsep=4pt]
  \item boundaries $-1,5$; True; \ans{In, and} $\Rightarrow$ \ans{$-1\le x\le 5$}\\ \gIn{-3}{7}{-1}{5}
  \item boundaries $-4,2$; False; \ans{Out, or} $\Rightarrow$ \ans{$x\le -4$ or $x\ge 2$}\\ \gOut{-6}{4}{-4}{2}
  \item boundaries $-3,2$; True; \ans{In, and} $\Rightarrow$ \ans{$-3< x< 2$} (open circles)\\ \gInOpen{-5}{4}{-3}{2}
  \item \ans{All real numbers}\\ \gAll{-4}{4}
\end{enumerate}
\end{multicols}

\end{tcolorbox}

\end{document}
